\section*{Bab \ref{ch:g10:functions} --- Fungsi}

\begin{solution}{exo:g10:functions:1}
$f(0) = 1$; \quad $f(2) = 8 - 6 + 1 = 3$; \quad
$f(-1) = 2 + 3 + 1 = 6$; \quad
$f\!\left(\tfrac12\right) = 2 \times \tfrac14 - \tfrac32 + 1 = \tfrac12 -
\tfrac32 + 1 = 0$.
\end{solution}

\begin{solution}{exo:g10:functions:2}
$f$: tanpa kendala, \omterm{def:g10:functions:function}{daerah asalnya} $\R$.

$g$: menuntut $x + 4 \neq 0$, \omterm{def:g10:functions:function}{daerah asalnya} semua \omterm{def:g10:numbers:sets}{bilangan real} kecuali $-4$.

$h$: menuntut $2x - 6 \geq 0$, yaitu $x \geq 3$: \omterm{def:g10:functions:function}{daerah asalnya}
$\intco{3}{+\infty}$.

$k$: $x^2 + 1 \geq 1 > 0$ tidak pernah lenyap, \omterm{def:g10:functions:function}{daerah asalnya} $\R$.
\end{solution}

\begin{solution}{exo:g10:functions:3}
\emph{\omterm{def:g10:functions:function}{Prapeta} dari $0$:} $x^2 - 2x = x(x-2) = 0$, jadi $0$ dan $2$.

\emph{\omterm{def:g10:functions:function}{Prapeta} dari $3$:} $x^2 - 2x = 3$, yaitu $x^2 - 2x - 3 = 0$, yaitu
$(x-3)(x+1) = 0$ (periksa dengan \omterm{def:g10:algebra:expand}{menjabarkannya}), jadi $3$ dan $-1$.

\emph{\omterm{def:g10:functions:function}{Prapeta} dari $-1$:} $x^2 - 2x + 1 = (x - 1)^2 = 0$, jadi \omterm{def:g10:functions:function}{prapeta}
tunggalnya $1$.
\end{solution}

\begin{solution}{exo:g10:functions:4}
$f(-1) = 1 + 2 + 1 = 4$: ya, $A(-1, 4)$ terletak pada \omterm{def:g10:functions:graph}{grafiknya}.
$f(2) = 4 - 4 + 1 = 1$: ya, $B(2, 1)$ pun terletak pada \omterm{def:g10:functions:graph}{grafiknya}.
\end{solution}

\begin{solution}{exo:g10:functions:5}
\emph{1.} Nilai terbesar dalam tabel itu adalah $4 = g(3)$ (\omterm{def:g10:functions:extrema}{maksimum}), yang
terkecil $-2 = g(0)$ (\omterm{def:g10:functions:extrema}{minimum}).

\emph{2.} Pada $\intcc{-3}{0}$ \omterm{def:g10:functions:function}{fungsi} itu \omterm{def:g10:functions:variations}{turun}, dan
$-1 < -0.5$, sehingga $g(-1) \geq g(-0.5)$.

\emph{3.} Pada setiap \omterm{def:g10:numbers:interval}{interval} kemonotonannya, $g$ mengambil setiap nilai
paling banyak sekali, sehingga $g(x) = 0$ mempunyai paling banyak satu
penyelesaian per \omterm{def:g10:numbers:interval}{interval}: seluruhnya paling banyak $3$
penyelesaian. (Di sini $0$ terletak antara $-2$ dan $1$, antara $-2$ dan
$4$, serta antara $0$ dan $4$, jadi ada tepat tiga.)
\end{solution}

\begin{solution}{exo:g10:functions:6}
Misalkan $u < v$. Maka
\[
f(v) - f(u) = (-2v + 7) - (-2u + 7) = -2(v - u) < 0,
\]
sebab $v - u > 0$. Jadi $f(v) < f(u)$: $f$ tegas \omterm{def:g10:functions:variations}{turun} pada $\R$.
\end{solution}

\begin{solution}{exo:g10:functions:7}
Lengkapkan kuadratnya: $(x+3)^2 = x^2 + 6x + 9$, jadi
\[
f(x) = x^2 + 6x + 11 = (x + 3)^2 + 2 .
\]
Untuk setiap $x$ berlaku $(x+3)^2 \geq 0$, jadi $f(x) \geq 2$; dan
$f(-3) = 0 + 2 = 2$. \omterm{def:g10:functions:extrema}{Minimum} $f$ adalah $2$, dicapai di $x = -3$.
\end{solution}

\begin{solution}{exo:g10:functions:8}
\emph{1.} Panjang $+$ lebar $= 20$, jadi panjangnya $20 - x$ dan
\[
A(x) = x(20 - x).
\]
Kedua ukurannya harus positif: $0 < x < 20$.

\emph{2.} $A(4) = 64$, $A(8) = 96$, $A(10) = 100$, $A(12) = 96$,
$A(16) = 64$. Nilainya \omterm{def:g10:functions:variations}{naik} sampai $x = 10$ lalu \omterm{def:g10:functions:variations}{turun} secara simetris:
luasnya tampak terbesar untuk $x = 10$, yaitu untuk kebun \emph{persegi}.
(\cref{ch:g11:quad} membuktikan dugaan ini.)
\end{solution}

\begin{solution}{exo:g10:functions:9}
\emph{1.} Untuk setiap $x$: $x^2 + 1 \geq 1 > 0$, jadi $f(x) = \frac{1}{x^2+1}$
bernilai positif; dan karena $x^2 + 1 \geq 1$, mengambil kebalikannya (kedua
ruas positif) memberi $f(x) \leq 1$.

\emph{2.} \omterm{def:g10:functions:extrema}{Maksimum}: $f(0) = 1$ dan $f(x) \leq 1$ untuk semua $x$, jadi $1$
\omterm{def:g10:functions:extrema}{maksimumnya}, dicapai di $0$. \omterm{def:g10:functions:extrema}{Minimum}: $f(x) > 0$ selalu, tetapi tidak ada
nilai $x$ yang mencapai $0$ (hasil bagi bilangan taknol bernilai taknol), dan
$f$ mengambil nilai sedekat yang kita mau ke $0$ untuk $x$ besar; jadi $f$
tidak mempunyai \omterm{def:g10:functions:extrema}{minimum}.
\end{solution}

\begin{solution}{exo:g10:functions:10}
\emph{1.} $f(v) - f(u) = v^2 - u^2 = (v-u)(v+u)$. Jika $0 \leq u < v$, maka
$v - u > 0$ dan $v + u > 0$ (jumlah bilangan taknegatif dan bilangan positif),
jadi $f(v) - f(u) > 0$: $f$ tegas \omterm{def:g10:functions:variations}{naik} pada $\intco{0}{+\infty}$.

\emph{2.} Jika $u < v \leq 0$, maka $v - u > 0$ tetap berlaku, tetapi kini
$v + u < 0$ (jumlah bilangan takpositif dan bilangan negatif), jadi
$f(v) - f(u) < 0$: $f$ tegas \omterm{def:g10:functions:variations}{turun} pada $\intoc{-\infty}{0}$.
\end{solution}

\begin{solution}{pb:g10:functions:1}
\textbf{1.} \omterm{def:g10:functions:function}{Daerah asal}: $x \neq 0$ (pembagian oleh $x$). $f(1) =
\frac12(1 + 2) = \frac32$; $f(2) = \frac12(2 + 1) = \frac32$;
$f\!\left(\frac32\right) = \frac12\left(\frac32 +
\frac43\right) = \frac{17}{12}$.

\textbf{2.} $\frac12\left(x + \frac2x\right) = \frac32$ memberi
$x^2 + 2 = 3x$, yaitu $x^2 - 3x + 2 = (x - 1)(x - 2) = 0$: jadi
\omterm{def:g10:functions:function}{prapeta} dari $\frac32$ adalah $1$ dan $2$ --- seperti sudah
diisyaratkan pertanyaan 1.

\textbf{3.} $f(x) = x$ memberi $x + \frac2x = 2x$, jadi
$\frac2x = x$, $x^2 = 2$: \omterm{pb:g10:functions:1}{titik tetapnya} $\sqrt2$ dan $-\sqrt2$.
Yang positif adalah diagonal persegi satuan, yang \omterm{ex:g10:numbers:classify}{irasional}
menurut soal akhir pekan tentang sifat \omterm{ex:g10:numbers:classify}{irasional} di jilid sebelumnya.

\textbf{4.} $1 \to \frac32 \to \frac{17}{12} \to
\frac{577}{408}$: \omterm{def:g10:numbers:approx}{hampiran} Heron bagi $\sqrt2$. Resepnya,
``rata-ratakan terkaan dengan $2/$terkaan'', tepat sama dengan
\emph{mengiterasi \omterm{def:g10:functions:function}{fungsi} $f$} --- dan bilangan yang dikejar
proses itu adalah \omterm{pb:g10:functions:1}{titik tetap} $f$.

\textbf{5.} \omterm{pb:g10:functions:1}{Titik tetapnya} adalah perpotongan \omterm{def:g10:functions:graph}{grafiknya}
dengan diagonal $y = x$: menyuapkan keluarannya kembali membuat
titik itu berjalan sepanjang \omterm{def:g10:functions:graph}{grafik} menuju perpotongan tadi. Pada
soal akhir pekan tentang dua termometer di jilid sebelumnya, angka
$-40^\circ$ adalah gagasan memotong diagonal yang sama, untuk \omterm{def:g10:functions:function}{fungsi}
konversinya.

\textbf{6.} $f(v) - f(u) = \frac{v - u}{2} + \frac1v - \frac1u
= \frac{v - u}{2} + \frac{u - v}{uv}
= (v - u)\left(\frac12 - \frac{1}{uv}\right)
= \frac{(v - u)(uv - 2)}{2uv}$. Untuk $0 < u < v \leq \sqrt2$:
$uv < 2$, isi kurungnya negatif, $f(v) < f(u)$: \omterm{def:g10:functions:variations}{turun}.
Untuk $\sqrt2 \leq u < v$: $uv > 2$: \omterm{def:g10:functions:variations}{naik}. Titik baliknya:
$\sqrt2$.

\textbf{7.} \omterm{def:g10:functions:variations}{Turun} sebelum $\sqrt2$, \omterm{def:g10:functions:variations}{naik} sesudahnya: $f$
mencapai \omterm{def:g10:functions:extrema}{minimumnya} di $\sqrt2$, bernilai
$f(\sqrt2) = \frac12\left(\sqrt2 + \frac{2}{\sqrt2}\right)
= \sqrt2$. Akibatnya: setiap keluaran mesin itu (untuk masukan
positif) paling sedikit $\sqrt2$ --- setelah satu putaran engkol,
semua terkaan berada di $\intco{\sqrt2}{+\infty}$.

\textbf{8.} Jika $x > \sqrt2$, maka $x^2 > 2$, sehingga
$\frac2x < x$ dan rata-rata $f(x)$ dari $x$ dan $\frac2x$
lebih kecil daripada $x$; lalu $f(x) \geq \sqrt2$ menurut pertanyaan 7,
dengan kesamaan hanya di $x = \sqrt2$. Jadi
$\sqrt2 < f(x) < x$: setiap iterasi memperbaiki secara tegas tanpa
kelewatan --- barisan pada pertanyaan 4 meluncur \omterm{def:g10:functions:variations}{turun} ke
$\sqrt2$.

\textbf{9.} Pada $\intoo{0}{+\infty}$: \omterm{def:g10:functions:variations}{turun} dari
$+\infty$ (dekat $0$) sampai \omterm{def:g10:functions:extrema}{minimum} $\sqrt2$ di
$x = \sqrt2$, lalu \omterm{def:g10:functions:variations}{naik} tanpa batas.

\textbf{10.} $g(x) = x$ memberi $x^2 = x$, $x(x - 1) = 0$: \omterm{pb:g10:functions:1}{titik
tetapnya} $0$ dan $1$. Dari $0.9$: $0.81$, $0.656$, $0.430$,
$0.185$ --- meluncur ke $0$. Dari $1.1$: $1.21$, $1.464$,
$2.144$, $4.595$ --- lari ke takhingga. \omterm{pb:g10:functions:1}{Titik tetap} $0$
menarik, $1$ menolak: beda sehelai rambut di awal, nasib yang
berlawanan.

\textbf{11.} $7 \to 22 \to 11 \to 34 \to 17 \to 52 \to 26 \to
13 \to 40 \to 20 \to 10 \to 5 \to 16 \to 8 \to 4 \to 2 \to 1$:
enam belas langkah, memuncak di $52$.

\textbf{12.} $27 \to 82 \to 41 \to 124 \to 62 \to 31 \to 94 \to
47 \to 142 \to 71 \to 214 \to \dots$ --- masih mendaki setelah
sepuluh langkah, dalam perjalanan menuju puncak $9\,232$ dan
penerbangan sepanjang $111$ langkah. Pelajarannya: aturan dua baris
dapat menghasilkan perilaku yang tak tertebak sekilas ---
bilangan bertetangga ($26$, $27$, $28$)
mempunyai penerbangan yang jauh berbeda.

\textbf{13.} Pangkat dari $2$ terbagi dua menuruni seluruh tangganya:
$2^k \to 2^{k-1} \to \dots \to 2 \to 1$. Di dasarnya,
$h(1) = 4$, $h(4) = 2$, $h(2) = 1$: penerbangannya masuk ke daur
$4 \to 2 \to 1 \to 4$. \omterm{pb:g10:functions:1}{Titik tetap} menuntut $\frac n2 = n$
(hanya $n = 0$) atau $3n + 1 = n$ (negatif): tidak ada di antara
\omterm{def:g10:numbers:sets}{bilangan bulat} positif --- penerbangan yang teramati berakhir bukan
pada \omterm{pb:g10:functions:1}{titik tetap}, melainkan pada \emph{daur} kecil itu.

\textbf{14.} Contoh penyangkalnya berupa penerbangan yang
membesar selamanya (tidak pernah kembali ke $1$) atau daur kedua
yang lepas dari $4 \to 2 \to 1$. Pemeriksaan sampai $10^{20}$
menyisakan tak hingga banyak bilangan yang belum teruji --- persis
pelajaran daerah lingkaran pada soal akhir pekan tentang ramalan
batang korek api di jilid sebelumnya: kecocokan pada berhingga
banyak kasus tidak membuktikan apa pun tentang semuanya.

\textbf{15.} Pada mesin Heron kita punya \emph{struktur} --- selisih
terfaktorkan (pertanyaan 6) yang menyingkap variasinya, sebuah
\omterm{def:g10:functions:extrema}{minimum}, sebuah apitan --- sedangkan aturan pergantian paritas mesin
hujan es tidak menawarkan pegangan semacam itu: belum ada yang
menemukan struktur yang menjinakkannya.

\textbf{16.} $\sqrt{x - 3}$: menuntut $x - 3 \geq 0$:
$\intco{3}{+\infty}$. $\frac{1}{x^2 - 9}$: menuntut
$x \neq \pm 3$: $\R$ tanpa $-3$ dan $3$.
$\sqrt{9 - x^2}$: menuntut $x^2 \leq 9$: $\intcc{-3}{3}$.

\textbf{17.} $f(x) = 3$ memberi $x + \frac2x = 6$, jadi
$x^2 - 6x + 2 = 0$, yaitu $(x - 3)^2 = 7$:
$x = 3 - \sqrt7$ atau $x = 3 + \sqrt7$ (keduanya positif: dua \omterm{def:g10:functions:function}{prapeta}
eksak).

\textbf{18.} Bermakna secara fisis selama batu itu melayang:
$H(t) = 5t(4 - t) \geq 0$ untuk $t \in \intcc{0}{4}$. Setelah dilengkapkan:
$H(t) = 20 - 5(t - 2)^2$: \omterm{def:g10:functions:extrema}{maksimum} $20$ m pada $t = 2$ s.
Variasinya: \omterm{def:g10:functions:variations}{naik} pada $\intcc{0}{2}$ dari $0$ ke $20$,
lalu \omterm{def:g10:functions:variations}{turun} pada $\intcc{2}{4}$ kembali ke $0$.

\textbf{19.} $20t - 5t^2 \geq 15 \iff t^2 - 4t + 3 \leq 0 \iff
(t - 1)(t - 3) \leq 0 \iff t \in \intcc{1}{3}$: dua detik penuh
di atas $15$ m.

\textbf{20.} Setiap mesin mempunyai \omterm{def:g10:functions:function}{daerah asal} (tempat aturannya
berlaku), sebuah aturan (rumus atau pergantian paritas), serta
\omterm{def:g10:functions:graph}{grafik} atau tabel yang menampilkannya utuh. Kepada masing-masing kita
ajukan dua pertanyaan besar: \emph{ke mana pengulangan bermuara}
(\omterm{pb:g10:functions:1}{titik tetap}, daur, tarikan --- terpecahkan untuk Heron, terbuka untuk
mesin hujan es) dan \emph{bagaimana keluaran bergerak mengikuti masukan}
(variasi, nilai ekstrem --- puncak batu, \omterm{def:g10:functions:extrema}{minimum} Heron).
Bab-bab berikutnya mempertajam keduanya: \omterm{def:g10:functions:function}{fungsi} acuan lebih dahulu,
lalu, pada tahun-tahun mendatang, turunan untuk mengukur variasi dan
limit untuk memastikan ke mana iterasi mendarat.

\end{solution}
